\documentclass[11pt]{article}

\usepackage[margin=1.12in]{geometry}
\usepackage{amsmath,amssymb,amsthm,mathtools}
\usepackage{microtype}
\usepackage[hidelinks]{hyperref}
\usepackage[nameinlink,noabbrev]{cleveref}

\newtheorem{theorem}{Theorem}[section]
\newtheorem{proposition}[theorem]{Proposition}
\newtheorem{lemma}[theorem]{Lemma}
\newtheorem{corollary}[theorem]{Corollary}
\newtheorem{definition}[theorem]{Definition}
\newtheorem{remark}[theorem]{Remark}

\newcommand{\C}{\mathbb C}
\newcommand{\R}{\mathbb R}
\newcommand{\N}{\mathbb N}
\newcommand{\HH}{\mathcal H}
\newcommand{\OO}{\mathcal O}
\newcommand{\WW}{\mathcal W}
\newcommand{\ord}{\operatorname{ord}}
\newcommand{\indm}{\operatorname{ind}_{-}}
\newcommand{\Res}{\operatorname{Res}}
\newcommand{\re}{\operatorname{Re}}
\newcommand{\im}{\operatorname{Im}}
\newcommand{\dd}{\,\mathrm d}

\title{\textbf{A Julia--Euler--Weil index law for off-critical Riemann divisors}}
\author{Author withheld for review}
\date{}

\begin{document}
\maketitle

\begin{abstract}
We determine the first normal orientation carried by a finite off-critical
divisor before any explicit-formula scalarization.  Given a finite packet of
genuinely complex functional-equation--Schwarz quartets
\[
 \tfrac12\pm a_j\pm i\gamma_j,\qquad
 0<a_j<\tfrac12,\quad \gamma_j>0,
\]
we form the Cayley phase of the strict-right pair and differentiate its
Julia defect in the unsquared aperture.  The derivative is pointwise
strictly positive.  Its integral, the pole price of the corresponding
rational Euler gauge, and the negative index of the isolated Weil packet
block satisfy
\[
 \frac{Q_{\mathrm{Julia}}}{4\pi}
   =\Delta\ord_{s=1}
   =\indm(\WW_{\mathrm{packet}})
   =2n.
\]
Thus the unsquared orientation exists, has a fixed sign, and is exactly
quantized; its polarity records an off-line defect rather than excluding
one.

On the subcone \(a_j^2+\gamma_j^2\leq1/4\), we additionally construct a
locally finite nonnegative continuous generalized-prime measure whose exact
prime-power transform generates the whole packet.  The associated
deformation retains the Riemann reflection law and has an entire
Schwarz-real hypertranscendental completion with precisely the inserted
zeros.  Hence Euler positivity, completed symmetry, hypertranscendence, and
positive first Julia orientation are jointly compatible with arbitrary
finite packets in this cone.  No criterion invariant under these positive
Euler surgeries can force the critical line.  What breaks the surgery is
not orientation but the exact simple-pole and discrete-prime normalization.
\end{abstract}

\section{Introduction}

The Riemann hypothesis is equivalent to the assertion that every
nontrivial zero of \(\zeta(s)\) has real part \(1/2\).  Three major
languages for this statement are usually kept separate.

First, the Euler side records primes through
\(-\zeta'/\zeta\).  Second, Weil's explicit formula turns zero localization
into positivity of a Hermitian quadratic form
\cite{Weil1952,Bombieri2000}.  Third, geometric current constructions
encode a zero divisor through Cauchy transforms or boundary Ward
responses.  Recent work on localized Weil forms, screw functions, and
spectral triples has made the relation between the first two languages
especially active \cite{ConnesConsani2021,CCM2025,Suzuki2026}.

A recurring difficulty is orientation.  Divisor currents and squared
normal coordinates retain the location and multiplicity of a zero but may
erase the sign of its transverse displacement from the critical line.
This suggests an apparently decisive repair: keep the normal coordinate
unsquared and prove a one-sided sign law before passing to a scalar
explicit formula.

The present paper answers that question at the first Julia derivative,
before divisor squaring or explicit-formula contraction.  The unsquared
aperture derivative exists and has a strict sign, but its sign is the sign
of the \emph{defect}, not a sign excluding the defect.  Every genuinely
off-critical quartet produces a positive normal derivative.  The same
quartet carries \(8\pi\) of integrated Julia--Ward charge, costs two units
of pole order at \(s=1\), and contributes a Weil block of negative index
two.  These are not three bounds.  They are three exact readings of one
integer.

The arithmetic realization is not attached after the index calculation.
On an explicit cone it starts with a nonnegative continuous
generalized-prime measure.  M\"obius inversion is performed at measure
level, the prime-power transform is recovered exactly, and its Laplace
transform is the logarithmic derivative of the rational divisor gauge.
The entire polynomial completion then compares Euler, geometric, and
quadratic readings of the same packet.  Outside that cone the index law
continues to hold, but positive measure realization is not asserted.

\subsection{Main result}

For \(a,\gamma>0\), put
\[
 q_{a,\gamma}(s)
 =\prod_{\epsilon,\eta\in\{\pm1\}}
   \left(s-\frac12-\epsilon a-i\eta\gamma\right),
 \qquad
 R_{a,\gamma}(s)
 =\frac{q_{a,\gamma}(s)}{[s(s-1)]^2}.
\]
For a finite nonempty set \(J\), let
\[
 R_J(s)=\prod_{j\in J}R_{a_j,\gamma_j}(s),
 \qquad
 F_J(s)=\zeta(s)R_J(s),
 \qquad
 X_J(s)=\xi(s)\prod_{j\in J}q_{a_j,\gamma_j}(s),
\]
where
\[
 \xi(s)=\frac12s(s-1)\pi^{-s/2}\Gamma(s/2)\zeta(s).
\]
The function \(X_J\) is the entire completion naturally associated with
\(F_J\).

Put \(c_j=a_j+\frac12\).  At boundary coordinate \(t\in\R\) and real
normal aperture \(\alpha\), define
\[
 \Phi_{j,t}(\alpha)
 =\prod_{\eta\in\{\pm1\}}
   \frac{c_j+i(t+\eta\gamma_j)-\alpha}
        {c_j+i(t+\eta\gamma_j)+\alpha},
 \qquad
 \Delta_{j,t}(\alpha)=1-|\Phi_{j,t}(\alpha)|^2.
\]
The variable \(\alpha\), not \(\alpha^2\), is the Julia normal coordinate.

\begin{theorem}[Finite-divisor Julia--Euler--Weil law]
\label{thm:main}
Let \(J\) be finite and nonempty, and suppose
\[
 0<a_j<\frac12,\qquad \gamma_j>0
\quad(j\in J).
\]
Then the following statements hold.

\begin{enumerate}
\item \(X_J\) is entire, satisfies
      \[
      X_J(1-s)=X_J(s),\qquad
      X_J(\overline{s})=\overline{X_J(s)},
      \]
      is hypertranscendental over \(\C(s)\), and vanishes at every
      \(\frac12\pm a_j\pm i\gamma_j\).
\item The first Julia normal derivative is
      \[
      \left.\frac{\partial}{\partial\alpha}
        \Delta_{j,t}(\alpha)\right|_{\alpha=0}
      =
      \frac{4c_j}{c_j^2+(t+\gamma_j)^2}
      +\frac{4c_j}{c_j^2+(t-\gamma_j)^2}>0.
      \]
      Its integral over \(t\in\R\) is \(8\pi\) for every \(j\), independent
      of height and horizontal displacement.
\item The rational gauge \(R_J\) has pole order \(2|J|\) at \(s=1\);
      \(F_J\) has total pole order \(1+2|J|\) there.
\item At every regular Cauchy probe \(z\) with
      \(\re z>0\) and \(\im z>0\), the polarized reciprocal-square
      response of the packet is strictly positive.
\item The isolated Weil packet block has signature
      \((2|J|,2|J|)\).
\item The integrated Julia charge, rational-gauge pole excess, and
      isolated Weil
      packet-block index obey
      \[
      \boxed{\quad
      \frac{Q_{\mathrm{Julia}}(J)}{4\pi}
       =\Delta\ord_1(F_J)
       =\indm(\WW_J)
       =2|J|.
      \quad}
      \]
\item If additionally
      \(a_j^2+\gamma_j^2\leq\frac14\) for every \(j\), there is a
      measurable locally integrable density \(p_J\geq0\) on
      \([0,\infty)\) whose generalized prime-power transform is exactly
      the logarithmic density attached to \(R_J\).  Away from poles and
      zeros,
      \[
      \frac{R_J'(s)}{R_J(s)}
      =-\sum_{j\in J}D_{a_j,\gamma_j}(s),
      \qquad
      D_{a,\gamma}(s)+D_{a,\gamma}(1-s)=0.
      \]
\end{enumerate}
\end{theorem}

The word ``isolated'' in the Weil clause is essential.  We compute the
quadratic block contributed by the inserted finite divisor.  We do not
assert, without additional separation hypotheses, that adding this block
to the unknown full zeta form increases its global negative index by
exactly \(2|J|\).  Bombieri's global finite-truncation result has its own
hypotheses \cite{Bombieri2000}.

\subsection{What the theorem decides}

The theorem rules out a specific inference:
\[
\begin{split}
&\text{positive Euler source}
+\text{Riemann reflection}
+\text{Schwarz reality}\\
&\qquad
+\text{entire hypertranscendental completion}
+\text{positive first Julia normal derivative}\\
&\hspace{45mm}\not\Longrightarrow
\text{zeros on the critical line}.
\end{split}
\]
This closes the possible loophole created by squaring the original normal
coordinate.  The pre-scalar, first-order sign is present and strict, but it
measures transverse divisor mass with the non-excluding polarity.

The conclusion is also narrower than a no-go theorem for the full
Connes--Consani--Moscovici or Hilbert--P\'olya programs.  Those programs
seek genuinely global operator limits and use the discreteness of the
ordinary primes.  The result here proves a no-go theorem for every
criterion invariant under the positive continuous Euler surgeries of
\cref{thm:main}.  It identifies exactly which data are surgery-blind and
which integer records the defect.  It does not claim that full CCM or
Hilbert--P\'olya constructions are surgery invariant.

\section{Positive generalized-prime realization}
\label{sec:euler}

\subsection{Prime powers and M\"obius inversion}

For a function \(p\) on \([0,\infty)\), define
\[
 (Tp)(u)=\sum_{n\geq1}\frac1{n^2}p(u/n).
\]
This is the logarithmic-coordinate prime-power transform.  Its formal
Dirichlet inverse is
\[
 (T^{-1}g)(u)
 =\sum_{n\geq1}\frac{\mu(n)}{n^2}g(u/n),
\]
where \(\mu\) is the ordinary M\"obius function.

The positivity of this inverse is not automatic.  The following
elementary criterion is useful.

\begin{lemma}[Monotone M\"obius positivity]
\label{lem:mobius}
If \(g:[0,\infty)\to[0,\infty)\) is nondecreasing, then
\[
 (T^{-1}g)(u)
 \geq\left(2-\frac{\pi^2}{6}\right)g(u)\geq0
 \qquad (u\geq0).
\]
Whenever the defining series are absolutely convergent,
\(T(T^{-1}g)=g\).
\end{lemma}

\begin{proof}
The \(n=1\) term is \(g(u)\).  For \(n\geq2\),
\(0\leq g(u/n)\leq g(u)\), and therefore
\[
 \left|
 \sum_{n\geq2}\frac{\mu(n)}{n^2}g(u/n)
 \right|
 \leq g(u)\sum_{n\geq2}\frac1{n^2}
 =\left(\frac{\pi^2}{6}-1\right)g(u).
\]
The lower bound follows.  Exact inversion is the identity
\[
 \sum_{d\mid n}\mu(d)=
 \begin{cases}
 1,&n=1,\\
 0,&n>1,
 \end{cases}
\]
after an absolutely convergent rearrangement.
\end{proof}

\subsection{The quartet density}

Put
\[
 \phi_{a,\gamma}(u)
 =\cosh(u/2)-\cosh(au)\cos(\gamma u)
 =\cosh(u/2)-\re\cosh((a+i\gamma)u).
\]
For \(u>0\), define the prime-power logarithmic density
\[
 g_{a,\gamma}(u)
 =\frac{4e^{u/2}}{u}\phi_{a,\gamma}(u),
 \qquad
 g_{a,\gamma}(0)=0.
\]

\begin{lemma}[Disc convexity]
\label{lem:convexity}
If \(a^2+\gamma^2\leq1/4\), then
\(\phi_{a,\gamma}\) is convex on \([0,\infty)\),
\(\phi_{a,\gamma}(0)=0\), and
\[
 u\longmapsto g_{a,\gamma}(u)
\]
is nonnegative and nondecreasing.
\end{lemma}

\begin{proof}
Let \(w=a+i\gamma\), so \(|w|\leq1/2\).  Differentiating twice gives
\[
 \phi_{a,\gamma}''(u)
 =\frac14\cosh(u/2)-\re\!\left(w^2\cosh(wu)\right).
\]
For \(u\geq0\),
\[
 \re\!\left(w^2\cosh(wu)\right)
 \leq |w|^2|\cosh(wu)|
 \leq\frac14\cosh(|w|u)
 \leq\frac14\cosh(u/2).
\]
Thus \(\phi''\geq0\).  Convexity and
\(\phi(0)=0\) imply that \(\phi(u)/u\) is nonnegative and nondecreasing
for \(u>0\).  Multiplication by the positive nondecreasing factor
\(4e^{u/2}\) proves the claim.
\end{proof}

Define the genuine prime density
\[
 p_{a,\gamma}(u)
 =\sum_{n\geq1}\frac{\mu(n)}{n^2}
     g_{a,\gamma}(u/n).
\]

\begin{proposition}[One-quartet source]
\label{prop:source}
Under the hypotheses of \cref{lem:convexity}, with \(a,\gamma>0\),
the function \(p_{a,\gamma}\) is measurable, nonnegative, and integrable
on every compact interval.  Moreover
\[
 Tp_{a,\gamma}=g_{a,\gamma}.
\]
\end{proposition}

\begin{proof}
Measurability follows termwise.  Nonnegativity is
\cref{lem:mobius}.  On a compact interval \([0,A]\), monotonicity gives
a uniform summable majorant proportional to
\(\sum n^{-2}g_{a,\gamma}(A)\), proving local integrability.
Absolute convergence also justifies M\"obius inversion.
\end{proof}

For the finite packet, set
\[
 p_J=\sum_{j\in J}p_{a_j,\gamma_j},
 \qquad
 g_J=\sum_{j\in J}g_{a_j,\gamma_j}.
\]
Finite addition preserves every conclusion of
\cref{prop:source}, and \(Tp_J=g_J\).

\subsection{Laplace transform and the rational gauge}

For \(\re s>1\), direct integration gives
\[
\begin{aligned}
D_{a,\gamma}(s)
&=\int_0^\infty e^{-su}\,u g_{a,\gamma}(u)\dd u\\
&=2\left(\frac1{s-1}+\frac1s\right)
 -\sum_{\epsilon,\eta\in\{\pm1\}}
  \frac1{s-\frac12-\epsilon a-i\eta\gamma}.
\end{aligned}
\]
The first line shows that \(D_{a,\gamma}\) is generated by a
nonnegative generalized von Mangoldt correction.  The second gives two
exact identities:
\[
 D_{a,\gamma}(s)+D_{a,\gamma}(1-s)=0,
 \qquad
 \frac{R_{a,\gamma}'(s)}{R_{a,\gamma}(s)}
 =-D_{a,\gamma}(s).
\]
Summing over \(J\) proves the source-to-gauge clauses of
\cref{thm:main}.

\begin{remark}[Relation to known Beurling counterexamples]
Positive continuous generalized-prime systems with off-line zeros are
known.  For example, Debruyne, Diamond, and Vindas use a continuous prime
measure whose zeta function is
\[
 \frac{(s-1)^2+1}{(s-1)^2},
\]
with zeros at \(1\pm i\) \cite[Example~4.3]{DebruyneDiamondVindas2018}.
The contribution here is not the bare existence of an off-line Beurling
zero.  It is the arbitrary finite FE--Schwarz packet realization, its
exact source-to-gauge identity, the first Julia orientation law, and the
common Euler--Ward--Weil index.
\end{remark}

\section{Completion and pole budget}
\label{sec:completion}

\subsection{Reflection, Schwarz symmetry, and hypertranscendence}

The root multiset of \(q_{a,\gamma}\) is invariant under
\(s\mapsto1-s\) and conjugation.  Hence
\[
 q_{a,\gamma}(1-s)=q_{a,\gamma}(s),\qquad
 q_{a,\gamma}(\overline{s})
 =\overline{q_{a,\gamma}(s)}.
\]
It follows immediately that \(X_J\) is entire and has the same two
symmetries.

\begin{proposition}[Prime-shift obstruction for the completed function]
\label{prop:xiht}
The completed Riemann function \(\xi\) is hypertranscendental over
\(\C(s)\).
\end{proposition}

\begin{proof}
Write \(L=\xi'/\xi\).  On \(\re s>1\), the Gamma recurrence and the Euler
series give the exact shift defect
\[
\begin{aligned}
L(s+2)-L(s)
&-\left(\frac1{s+2}+\frac1{s+1}-\frac1{s-1}\right)\\
&=\sum_{n\geq1}
 \frac{\Lambda(n)(1-n^{-2})}{n^s}
 =:V(s).
\end{aligned}
\]
The coefficient of \(p^{-s}\) is
\(\log p\,(1-p^{-2})\neq0\) at every prime \(p\), while the support is
contained in the prime powers.  Ostrowski's Dirichlet-series theorem
therefore implies that \(V\) is hypertranscendental over \(\C(s)\)
\cite{Ostrowski1920}; modern difference--differential versions apply as
well \cite{Nagoshi2016}.

If \(\xi\) were differentially algebraic over \(\C(s)\), then so would be
\(L=\xi'/\xi\).  Translation \(s\mapsto s+2\) preserves differential
algebraicity over \(\C(s)\), and the displayed rational correction is
differentially algebraic.  The left side would make \(V\)
differentially algebraic, a contradiction.
\end{proof}

Multiplication by a nonzero polynomial preserves hypertranscendence:
if \(P\xi\) satisfied an algebraic differential equation over \(\C(s)\),
the identities obtained by differentiating \(\xi=(P\xi)/P\) would make
\(\xi\) differentially algebraic.  Consequently \(X_J\) is
hypertranscendental.

This point matters for route selection.  The completed function is
hypertranscendental because a two-step translation exposes infinitely many
prime frequencies.  Yet the same obstruction survives every finite
off-critical packet in \cref{thm:main}.  Differential-algebraic complexity
therefore detects the infinite prime source while remaining blind to
finite horizontal divisor surgery.

\subsection{Exact pole order}

At \(s=1\),
\[
 q_{a,\gamma}(1)
 =\left(\left(\tfrac12-a\right)^2+\gamma^2\right)
  \left(\left(\tfrac12+a\right)^2+\gamma^2\right)>0.
\]
Therefore each \(R_{a,\gamma}\) has a pole of order exactly two at
\(s=1\).  Products add meromorphic orders:
\[
 \ord_{1}R_J=-2|J|.
\]
Since \(\zeta\) has a simple pole at \(1\),
\[
 \ord_{1}F_J=-(1+2|J|).
\]
We call
\[
 \Delta\ord_1(F_J)=2|J|
\]
the pole-order excess of the packet.

The completion identity is
\[
 \frac12[s(s-1)]^{2|J|+1}
   \pi^{-s/2}\Gamma(s/2)F_J(s)
 =X_J(s).
\]
Thus the higher pole is not an artefact hidden outside the completion:
it is exactly the linear budget paid to insert the finite divisor while
retaining positivity and reflection.

\section{Reciprocal-square Cauchy polarization}
\label{sec:orientation}

Use the centered spectral coordinate
\[
 A=\gamma-ia
\]
for the upper-right zero \(\frac12+a+i\gamma\).  The critical line is
the real axis in this coordinate.  For a Cauchy probe \(z\in\C\), define
the contribution of a spectral point \(A\) by
\[
 \tau_z(A)
 =-\frac{\im(A^{-2})}{|A-z|^2}
 =\frac{2(\re A)(\im A)}
 {|A|^4|A-z|^2}.
\]
Although this reader uses the reciprocal square, taking its imaginary
polarization retains a signed transverse factor.  It is a secondary check
on whether the later Cauchy scalarization reverses the first Julia
orientation.

For the orbit
\(\{A,\overline A,-A,-\overline A\}\), put
\[
\begin{aligned}
V_z(A)
={}&
\frac{(\re A)^2(\im A)^2(|A|^2+|z|^2)}
{|A|^4
 |A-z|^2|\overline A-z|^2
 |A+z|^2|\overline A+z|^2}.
\end{aligned}
\]
At regular probes all denominators are nonzero, and \(V_z(A)\geq0\).

\begin{proposition}[Universal quartet orientation identity]
\label{prop:orientation}
For every regular \(A,z\),
\[
 \tau_z(A)+\tau_z(\overline A)+\tau_z(-A)+\tau_z(-\overline A)
 =64(\re z)(\im z)V_z(A).
\]
Moreover \(V_z(A)=0\) if and only if
\(\re A=0\) or \(\im A=0\).
\end{proposition}

\begin{proof}
Substituting the displayed rational formula for \(\tau_z\), putting the
four terms over a common denominator, and collecting the numerator gives
the identity.  Every denominator factor is strictly positive at a regular
probe.  The numerator of \(V_z(A)\) is
\[
 (\re A)^2(\im A)^2(|A|^2+|z|^2).
\]
Since \(A\neq0\), the vanishing statement follows.
\end{proof}

For \(A_j=\gamma_j-ia_j\), both coordinates are nonzero.  If
\(\re z,\im z>0\), \cref{prop:orientation} gives
\[
 \OO_J(z)
 =64(\re z)(\im z)\sum_{j\in J}V_z(A_j)>0.
\]
Thus the reciprocal-square polarization has the same positive defect
polarity as the first Julia derivative below.  Its strict positivity is
unconditional, but it detects the presence of an off-critical quartet.
A zero-exclusion argument based on a nonpositive response would require
the opposite inequality \(\OO_J(z)\leq0\).  The source-positive examples
show that such an inequality cannot follow from the preserved data.

\section{The first Julia derivative and Ward charge}
\label{sec:ward}

For \(c=a+\frac12>0\), the squared modulus of the phase introduced before
\cref{thm:main} is
\[
 |\Phi_t(\alpha)|^2
 =\prod_{\eta\in\{\pm1\}}
 \frac{(c-\alpha)^2+(t+\eta\gamma)^2}
      {(c+\alpha)^2+(t+\eta\gamma)^2}.
\]
For every \(\alpha\geq0\), each Cayley factor is contractive.  Direct
differentiation at the undeformed face gives
\[
\begin{aligned}
\left.\frac{\partial}{\partial\alpha}
 (1-|\Phi_t(\alpha)|^2)\right|_{\alpha=0}
 &=
 \frac{4c}{c^2+(t+\gamma)^2}
 +\frac{4c}{c^2+(t-\gamma)^2}.
\end{aligned}
\]
This is the first derivative in the normal aperture.  No squaring of
\(\alpha\), zero divisor, or explicit-formula test function enters the
calculation.

The strict-right upper Schwarz pair of one quartet has the edge response
\[
\mathcal R_{a,\gamma}(w)
=2\left(
\frac1{w+a+i\gamma}
+\frac1{w+a-i\gamma}
\right).
\]
On the Euler edge \(w=\frac12+it\), define the infinitesimal Ward density
\[
\begin{aligned}
\omega_{a,\gamma}(t)
=2\bigg[
&\frac{2(a+\frac12)}
 {(a+\frac12)^2+(t+\gamma)^2}\\
&+\frac{2(a+\frac12)}
 {(a+\frac12)^2+(t-\gamma)^2}
\bigg].
\end{aligned}
\]
It is twice the real part of \(\mathcal R_{a,\gamma}\) on that edge and
is pointwise positive.  The preceding calculation proves that it is
exactly the Julia phase-defect derivative:
\[
 \omega_{a,\gamma}(t)
 =\left.\partial_\alpha
   (1-|\Phi_t(\alpha)|^2)\right|_{\alpha=0}.
\]

\begin{proposition}[Ward quantization]
\label{prop:ward}
For every \(a>-\frac12\) and every \(\gamma\in\R\),
\[
 \int_{-\infty}^{\infty}\omega_{a,\gamma}(t)\dd t=8\pi.
\]
Consequently
\[
 Q_{\mathrm{Julia}}(J)
 :=\sum_{j\in J}\int_{\R}
 \left.\partial_\alpha
   (1-|\Phi_{j,t}(\alpha)|^2)\right|_{\alpha=0}\dd t
 =8\pi|J|.
\]
\end{proposition}

\begin{proof}
Each shifted Lorentzian has integral
\[
 \int_{-\infty}^{\infty}
 \frac{2c}{c^2+(t-b)^2}\dd t=2\pi
 \qquad(c>0).
\]
There are two such terms and one outer factor \(2\).
\end{proof}

The charge is independent of height and horizontal displacement.  It
counts the orbit, but by itself supplies no upper bound and no vanishing
theorem.  Crucially, the failure occurs already at the unsquared first
normal derivative: later scalarization does not create the wrong polarity.
Its significance emerges after comparison with pole budget and Weil block.

\section{The isolated Weil packet block}
\label{sec:weil}

For centered divisor coordinate \(w=\rho-\frac12\), the zero side of
Weil's form contains products of transform values at \(w\) and \(-w\).
The finite-dimensional algebra attached to one symmetry orbit has three
types.

\begin{itemize}
\item A critical pair \(\{i\gamma,-i\gamma\}\) gives
      \(2(X^2+Y^2)\), which is positive semidefinite.
\item A real pair \(\{a,-a\}\) gives \(2XY\), of signature \((1,1)\).
\item A generic quartet
      \(\{w,-w,\overline w,-\overline w\}\) gives
      \[
      4(Z_1W_1-Z_2W_2),
      \]
      of signature \((2,2)\).
\end{itemize}

The last assertion is transparent from
\[
\begin{aligned}
4(Z_1W_1-Z_2W_2)
={}&(Z_1+W_1)^2-(Z_1-W_1)^2\\
&-(Z_2+W_2)^2+(Z_2-W_2)^2.
\end{aligned}
\]
It is nonnegative on the two-plane
\((W_1,W_2)=(Z_1,-Z_2)\) and negative definite on
\((W_1,W_2)=(-Z_1,Z_2)\).

Thus one generic quartet has isolated packet-block negative index two.
For \(|J|\) independent orbit blocks,
\[
 \WW_J=\bigoplus_{j\in J}\WW_j,
 \qquad
 \operatorname{signature}(\WW_J)=(2|J|,2|J|),
 \qquad
 \indm(\WW_J)=2|J|.
\]

This is the local algebra behind Bombieri's theorem that, if RH has only
finitely many exceptions, sufficiently large truncations of the global
Weil form have one negative eigenvalue for every two exceptional zeros
\cite{Bombieri2000}.  Our use of it is deliberately packetwise and does
not import a global index-additivity assertion.

\section{The Julia--Euler--Weil law}
\label{sec:index}

Combining \cref{sec:completion,sec:ward,sec:weil} gives the central
identity.

\begin{theorem}[Julia--Euler--Weil packet index]
\label{thm:index}
For every finite packet in \cref{thm:main},
\[
\frac{Q_{\mathrm{Julia}}(J)}{4\pi}
=\Delta\ord_1(F_J)
=\indm(\WW_J)
=2|J|.
\]
\end{theorem}

\begin{proof}
\Cref{prop:ward} gives
\(Q_{\mathrm{Julia}}/(4\pi)=2|J|\).
The exact order calculation gives
\(\Delta\ord_1=2|J|\).
The generic-orbit diagonalization gives
\(\indm(\WW_J)=2|J|\).
\end{proof}

The identity is a finite-divisor index theorem.  Its three entries have
different origins:
\[
\begin{array}{c|c|c}
\text{language}&\text{object}&\text{integer}\\ \hline
\text{Euler source}&\text{pole at }s=1&2|J|\\
\text{Julia boundary}&\text{normal defect current}&Q/(4\pi)=2|J|\\
\text{explicit formula}&\text{generic Weil blocks}&\indm=2|J|.
\end{array}
\]
The polarized Cauchy response in \cref{sec:orientation} supplies a fourth,
non-quantized reading: it is a strictly positive weighted version of the
same orbit census.

\section{A surgery obstruction for RH criteria}
\label{sec:nogo}

Let a \emph{positive Euler quartet surgery} mean multiplication of a
zeta function by one of the rational factors \(R_{a,\gamma}\) under the
disc condition \(a^2+\gamma^2\leq1/4\), together with multiplication of
its entire completion by \(q_{a,\gamma}\).

\begin{definition}
A predicate \(\mathcal P\) on a generalized Euler source and its completed
function is \emph{packet-surgery invariant} if
\(\mathcal P\) remains true after every finite positive Euler quartet
surgery.
\end{definition}

\begin{corollary}[Surgery no-go]
\label{cor:nogo}
No packet-surgery-invariant predicate that holds for the Riemann source
and completion can imply that all zeros of the completion lie on the
critical line.
\end{corollary}

\begin{proof}
Apply one nonempty surgery from \cref{thm:main}.  Invariance preserves
\(\mathcal P\), while the new completion has the four prescribed
off-critical zeros.
\end{proof}

The force of the corollary depends on which properties are invariant or
survive the surgery.  The theorem proves this simultaneously for:
\[
\begin{gathered}
\text{nonnegative locally finite continuous generalized-prime source},\\
\text{exact source-to-logarithmic-derivative identity},\\
\text{Riemann reflection and Schwarz reality},\qquad
\text{entire completion},\\
\text{hypertranscendence over }\C(s),\qquad
\text{strictly positive first Julia normal derivative}.
\end{gathered}
\]
Accordingly, no conjunction of these data can supply zero exclusion.

Pole order is not invariant: every generic quartet costs exactly two
units.  This is the sharp boundary exposed by the index law.  Within this
finite positive-surgery class, retaining a simple pole at \(s=1\) forbids
every nonempty insertion.  That statement is a rigidity theorem for the
deformation class, not a proof of RH for the undeformed zeta function.
It does, however, locate the missing information precisely: a successful
source argument must use a feature not stable under continuous positive
Euler surgery, such as the exact simple-pole normalization coupled to the
discrete atomic arithmetic source.

\section{Relation to current RH programs}

Weil's criterion is global: positivity for the full test-function class is
equivalent to RH \cite{Weil1952}.  Bombieri's variational and index
analysis makes the negative directions caused by finitely many exceptions
quantitative \cite{Bombieri2000,Bombieri2003}.  Suzuki's screw-function
framework gives a continuous-kernel realization of localized Weil forms
and relates them to self-adjoint extensions \cite{Suzuki2026}.
Connes, Consani, and Moscovici construct finite-prime self-adjoint
operators whose observed spectra approximate low zeta zeros; proving the
required global convergence would establish RH \cite{CCM2025}.

\Cref{thm:index} does not contradict those programs.  It separates two
logical levels.

\begin{enumerate}
\item At finite divisor level, Julia--Ward charge, Euler pole budget, and the
      generic Weil signature are the same defect integer.  Re-encoding one
      as another cannot create a vanishing theorem.
\item A global spectral proof must control a limit that is not invariant
      under the continuous positive Euler surgeries constructed here.
      In particular, it must retain arithmetic information beyond
      reflection, positivity, hypertranscendence, and finite divisor
      orientation.
\end{enumerate}

Connes and Consani already show that a tempting placewise sign inequality
does not hold in general and emphasize the role of global Poisson
normalization \cite{ConnesConsani2021}.  The present result gives a
complementary divisor-level mechanism: the one-sided normal sign exists,
but positive Euler sources realize it with the off-critical polarity.

\section{Scope and conclusion}

Four scope boundaries are part of the theorem.

\begin{enumerate}
\item The index and first-Julia-orientation laws hold for every finite
      generic quartet packet in the open critical strip.  Finiteness is
      essential to the present theorem.
\item The nonnegative generalized-prime realization is continuous, not a
      Beurling prime sequence, and is asserted on the centered closed disc
      \(a^2+\gamma^2\leq1/4\), where M\"obius inversion is positive by the
      monotonicity criterion.
\item The Weil index is the exact signature of the isolated inserted
      packet block, not an unconditional global index difference for the
      unknown zeta divisor.
\item The theorem is a no-go for packet-surgery-invariant criteria, not
      for every possible CCM, trace-formula, or Hilbert--P\'olya
      construction.
\end{enumerate}

Within those boundaries, the orientation problem is closed.  Squaring the
normal coordinate can erase transverse sign, but restoring the unsquared
first derivative does not force the critical line.  Every source-positive
off-critical quartet has a strict positive Julia response before
scalarization.  Its geometric charge, Euler pole price, and Weil block
index are locked by
\[
 \frac{Q_{\mathrm{Julia}}}{4\pi}
 =\Delta\ord_1
 =\indm(\WW_{\mathrm{packet}}).
\]
The equality identifies a universal finite-surgery defect invariant.
Any RH proof operating in these languages must introduce a selector that
breaks positive Euler quartet surgery; merely preserving more faithfully
the orientation already present cannot suffice.

\begin{thebibliography}{99}

\bibitem{Bombieri2000}
E. Bombieri,
\newblock Remarks on Weil's quadratic functional in the theory of prime
numbers, I,
\newblock \emph{Rend. Lincei Mat. Appl.} \textbf{11} (2000), 183--233.

\bibitem{Bombieri2003}
E. Bombieri,
\newblock A variational approach to the explicit formula,
\newblock \emph{Comm. Pure Appl. Math.} \textbf{56} (2003), 1151--1164.

\bibitem{CCM2025}
A. Connes, C. Consani, and H. Moscovici,
\newblock Zeta spectral triples,
\newblock arXiv:2511.22755 (2025).

\bibitem{ConnesConsani2021}
A. Connes and C. Consani,
\newblock Weil positivity and trace formula, the archimedean place,
\newblock \emph{Selecta Math.} \textbf{27} (2021), Article 77.

\bibitem{DebruyneDiamondVindas2018}
G. Debruyne, H. G. Diamond, and J. Vindas,
\newblock \(M(x)=o(x)\) estimates for Beurling numbers,
\newblock \emph{J. Th\'eor. Nombres Bordeaux} \textbf{30} (2018),
469--483.

\bibitem{HilberdinkLapidus2006}
T. W. Hilberdink and M. L. Lapidus,
\newblock Beurling zeta functions, generalised primes, and fractal
membranes,
\newblock \emph{Acta Appl. Math.} \textbf{94} (2006), 21--48.

\bibitem{Nagoshi2016}
H. Nagoshi,
\newblock Hypertranscendence of \(L\)-functions for
\(\mathrm{GL}_m(\mathbb A_{\mathbb Q})\),
\newblock \emph{Bull. Aust. Math. Soc.} \textbf{93} (2016), 388--399.

\bibitem{Ostrowski1920}
A. Ostrowski,
\newblock \"Uber Dirichletsche Reihen und algebraische
Differentialgleichungen,
\newblock \emph{Math. Z.} \textbf{8} (1920), 241--298.

\bibitem{Suzuki2026}
M. Suzuki,
\newblock Weil's quadratic form via the screw function,
\newblock arXiv:2606.09096 (2026).

\bibitem{Weil1952}
A. Weil,
\newblock Sur les ``formules explicites'' de la th\'eorie des nombres
premiers,
\newblock \emph{Medd. Lunds Univ. Mat. Sem.} (1952), 252--265.

\end{thebibliography}

\end{document}
